\section{Introduction}\label{sec:introduction}

\sumpara A planner on a flow line must commit to one job order before processing times are known, and the common dispatching rule ignores most of the information the planner holds. \sumparaend
\fin On flow lines such as printed circuit board assembly or furniture finishing, every job visits the same stages in the same order, and when jobs cannot pass one another, the order chosen at the first stage holds at every stage. \fine
\fin A job that is quick at the first stage and slow at a later one then delays every job behind it, and the resulting waiting accumulates downstream. \fine
\fin The planner of such a line fixes one sequence for the jobs of a planning period before processing begins, although the jobs are released at different times, carry due dates of unequal importance, and take processing times that will differ from their planned values. \fine
\fin Dispatching rules are used very often in practice, and the earliest due date (EDD) rule is among the most common rules for tardiness objectives~\cite{vallada2008review}. \fine
\fin EDD orders the jobs by due date alone, so it cannot account for job weights, release dates, or processing times at the stages after the first. \fine

\sumpara Searched sequences are known to beat dispatching rules on exact data, but whether that advantage survives variable processing times has been studied only for the makespan. \sumparaend
\fin Minimizing total tardiness in a permutation flow shop is NP-hard in the ordinary sense with one machine and in the strong sense with two or more~\cite{vallada2008review}, and the iterated greedy algorithm of \citet{ruiz2008igwt} was the best of the methods its authors tested for weighted tardiness. \fine
\fin That evaluation, like the review of \citet{vallada2008review}, used instances with fixed processing times and no release dates. \fine
\fin Studies of flow shops with random processing times have concentrated on the expected makespan, for which heuristics adapted from deterministic procedures, including sequences computed from mean processing times, perform close to the best stochastic methods~\cite{baker2012stochastic,framinan2015stochastic,juan2014simheuristic}, and a review of 100 studies of flow shops under uncertainty identifies tardiness and its weighted form as criteria that remain to be examined~\cite{gonzalezneira2017review}. \fine
\fin A broader review of schedule execution concludes that optimized predictive schedules outperform simple rules when uncertainty is low and lose that advantage beyond a level of variability that depends on the system~\cite{aytug2005executing}. \fine
\fin For weighted tardiness with release dates, it therefore remains open whether a searched sequence, held fixed, keeps its advantage over the EDD sequence as processing times vary, and at which level of variability that advantage weakens. \fine

\sumpara We compare the two fixed sequences on exact data, certify small optima, and execute both on paired scenarios across a range of variability. \sumparaend
\fin We adapt the iterated greedy algorithm of \citet{ruiz2008igwt} to an eight-stage permutation flow shop with release dates and total weighted tardiness, and we compare it with EDD on four synthetic instances with 6 to 120~jobs. \fine
\fin A mixed-integer program solved with HiGHS~\cite{highs} certifies the optimal values of the two smallest instances, and a checker that shares no code with either method recomputes the feasibility and the measures of every schedule. \fine
\fin We then execute both sequences, unchanged, on 180 paired scenarios of the 45-job instance in which each processing time is multiplied by an independent lognormal factor with mean one, at coefficients of variation from 0.1 to 0.8. \fine

\sumpara The searched sequence nearly eliminates the rule's tardiness on exact data and keeps its mean advantage at every level of variability tested, while its advantage in the worst periods erodes. \sumparaend
\fin On exact data, iterated greedy reduces total weighted tardiness relative to EDD by 98.5, 97.6, and 81.0~percent on the instances with 15, 45, and 120~jobs, and it attains the proven optima on the instances with 6 and 15~jobs. \fine
\fin Across the 180 paired scenarios, the iterated greedy sequence has the lower weighted tardiness in 179, and its mean reduction stays between 2,572.0 and 3,683.7 weighted minutes at every coefficient of variation tested. \fine
\fin The relative reduction, however, falls from 94.6~percent at a coefficient of variation of 0.1 to 27.2~percent at 0.8, and at 0.8 the 10th percentile of the paired reduction falls to 851.5 weighted minutes. \fine

\sumpara The contribution is a mean-versus-tail account of how variability changes the value of a searched fixed sequence for weighted tardiness. \sumparaend
\fin The results extend the finding for the expected makespan, that sequences computed on mean processing times remain effective, to total weighted tardiness with release dates, and they locate the erosion that \citet{aytug2005executing} describe in the tail of the per-period gain before its mean. \fine
\fin For a planner, the searched sequence is worth adopting at every level of variability we tested, and the variability of the line is the quantity that determines how dependably its gain is realized in a given period. \fine
